% This section lists a number of 2-3 research questions that the thesis should address. 
% The research questions should be put into the context of the idea, and you should explain 
% what they mean and how you want to address each. We also suggest including an example scenario 
% of a successful outcome of the thesis, i.e., what would your tool allow us to do once it exists?

\chapter{Research questions} \label{research_questions}


This paper will address the following research questions:
\begin{enumerate}
    \item How can automated model checking and other formal-verification 
        techniques be used for verifying liveness (replication reliability),
        safety and other properties in self-replicating robotic systems(SRRS), 
        and what domain-specific adaptations are required for practical implementation?
    \item What systematic methodology can translate the behaviors and 
        structural constraints of self-replicating robots into expressive 
        temporal-logic models and specifications?
    \item What fundamental limitations—computational, semantic, or 
        modelling-related—do current face when applied to highly dynamic, 
        self-replicating robotic collectives, and how might these limitations be mitigated?
\end{enumerate}


\section{Analysis of model checking approaches in self-replicating systems domain}

% The first question includes examination of the state-of-the-art 
% of model checking approaches to determine whether they can be used to model and verify 
% an SRR system.
%
% The experience and different practical approaches in the field of self-reproducing 
% automata and robots should be studied, based on which the main 
% properties that should be included in the specifications have to be identified.
%
% Initially can be considered basic properties like liveness (replication reliability) and safety.
% With some adaptations it should be possible to express it using every model-checking
% specification language.
% During this analysis or further experimental checking, 
% more complex properties can be discovered  and integrated.
The first research question involves a comprehensive survey and critical 
evaluation of current state-of-the-art model checking and formal verification 
tools, such as SPIN, NuSMV, and PRISM, focusing on their applicability to modeling 
and verifying the unique dynamics of self-replicating robotic systems (SRRS). 

This investigation includes a review of relevant literature and existing systems 
in self-replicating robotics and cellular automata to identify key operational 
characteristics and formalizable behaviors. It also involves identifying 
core properties relevant to SRRS, such as liveness (ensuring successful 
replication within time or resource constraints), safety (avoiding 
resource conflicts or hazardous outcomes), as well as broader system-level guarantees 
like termination, deadlock-freedom, and fault-tolerance. 
The study will further explore how these properties can be 
expressed in various formal languages, including Linear 
Temporal Logic (LTL), Computation Tree Logic (CTL), and probabilistic extensions. 
This will require examining what adaptations or abstractions are needed to 
make these languages suitable for SRRS modeling. 

Iterative experimentation with selected tools will help assess the feasibility 
of encoding and verifying these properties and may lead to the discovery 
of more complex or emergent behaviors to integrate in later stages.

\section{Experimental checking}



The second research question aims to define a systematic formal modeling 
and verification workflow tailored to SRRS. 
This involves developing a formal representation of 
self-replication processes, structural rules, environmental interactions, 
and resource usage. 

It also includes establishing a methodology to map real-world 
or simulated robotic behaviors - such as perception, manipulation, and actuation- onto 
abstract states and transitions that are amenable to formal analysis. 
Based on this framework, appropriate model-checking tools and specification 
languages will be selected and justified through experimental evaluation within 
simulation environments. 

A prototype or proof-of-concept implementation will be developed 
using a simplified SRRS scenario, such modular robotic replication, 
to demonstrate how this methodology captures essential properties and enables formal verification. 

The ultimate goal is to bridge the gap between low-level robotic simulation and control 
systems and high-level formal models, ensuring that the designed 
behaviors are both physically feasible and logically justified.


\section{Limitations and challenges}


The third research question addresses the fundamental limitations that 
current formal methods and model-checking tools face when applied to the 
complex domain of self-replicating robotic systems.
It includes analyzing computational constraints such as state-space explosion, 
undecidability in dynamic conditions, and potential inefficiencies in scaling existing verification tools. 

The study will also examine semantic mismatches between physical 
or simulated robot behavior and abstract models, particularly in representing kinematic 
and dynamic aspects, sensor errors, and asynchronous or partially 
observable environments. 

Additionally, the research will explore modeling limitations related 
to the representation of dynamic structure generation and adaptive behaviors.

To address these challenges, various mitigation strategies will be proposed,
including modular verification, changes in system abstractions  
and hybrid techniques that combine formal models with simulation-based validation. 

The outcome of this part may include recommendations for extending the 
selected verification toolsets, adapting implementation strategies, or designing 
new experimental setups specifically for the analysis of SRRS.



